%% Research design sections 6 and 7 -- elimination that is safe, ranking that
%% is only ranking, and an adaptive K.
%%
%% Contributions B and C. The rule that binds this section: a pruning stage is
%% a RELAXATION of the feasibility test and may reject only when the most
%% optimistic arithmetic already misses a declared constraint. The throughput
%% bound exists only because H1 added slo.min_goodput_rps and must not run at
%% all when that field is unset.
%%
%% The five states never merge, and the section says why: impossible_proven,
%% excluded_by_scope, deferred_heuristic, unknown_measurement, evaluated. A
%% budget is a property of the search, not of the hardware.
%%
%% GS-12 goes here as "the initial hypothesis was ..., and the measurement said
%% otherwise": the ranker's goodput term divided by the bound's optimistic
%% ceiling, which admitted four false positives per fixture.
%%
%% Numbers: tables/e_g1b_topk.tex, tables/e_g2_topk_holdout.tex,
%% tables/e_g7_holdout.tex.
\section{Provable elimination and an adaptive budget}
\label{sec:search}

% claim: C2
\paragraph{Elimination is a relaxation, or it is not elimination.}
A candidate may be removed only when arithmetic that is optimistic on every
uncertain term still misses a constraint the specification declares. Five
checks qualify: backend compatibility, device memory, communication latency
against the time-to-first-token target, throughput capacity against a declared
goodput floor, and a cost lower bound against the incumbent. Each publishes the
relaxations it made---``every KV slot holds a live request'', ``a decode step
costs only the weight sweep'', ``cut capacity is nominal and uncontended''---so
a reader can check that the arithmetic really is optimistic rather than merely
plausible.

Anything that cannot be shown to be a relaxation is not called proved. It is
reported as \texttt{deferred\_heuristic}, in a different column, and the
candidates it removes are not counted among those shown infeasible. This is a
rule inherited from the planner being extended, which had once removed
under-provisioned candidates with a throughput bound its specification did not
declare; the bound here runs only when the specification sets a goodput floor,
and does not run at all when that field is unset.

\paragraph{The cut bound.}
The graph makes one check possible that an island-level planner cannot express:
the maximum flow between a candidate's device sets, over nominal capacities
minus external reservations, is a ceiling on what its traffic can achieve. A
candidate whose demand exceeds that ceiling misses its target under the most
favourable routing available. That is a proof relative to the declared cluster
description---its capacities and its reservations---and not an estimate; it is
as right about the hardware as the description is.

% claim: C7
% claim: C8
\paragraph{Ranking is only ranking.}
What survives elimination is ordered by a service-margin estimate: for each
SLO, the ratio of what the candidate is estimated to need to what the
specification allows, and the ranking reads the maximum of those ratios---the
binding constraint---because a candidate meets its objectives only if every one
is met. The ranker produces no verdict; a candidate it places last is not a
candidate it rejects.

The estimate was wrong once, in a way worth recording. Its goodput term divided
by the elimination bound's optimistic ceiling rather than by a knob-aware
throughput estimate, which made every candidate look comfortable on that axis
and admitted false positives on each diagnosis fixture. The correction was
diagnosed on two fixtures and then checked on two it had never seen, because a
correction validated on the data that motivated it has not been validated.

% claim: C9
Neither correctness number moves when the ranker changes, and that is asserted
by a test rather than argued: ranking decides what is evaluated first, not what
is eliminated, so a ranker change that moved \texttt{false\_infeasible} would
be evidence of a different bug.

% claim: C5
% claim: C6
\paragraph{An adaptive budget.}
The search evaluates representatives in ranked order under a schedule of
increasing $K$, stopping when the budget is exhausted or when no unevaluated
candidate could beat the incumbent. It reports which of those happened.

Against the existing planner's surrogate top-$K$ on a specification whose SLOs
bind, this reaches higher feasible recall at every $K$ on the toy corpus (the
diagnosis fixtures the correction was made on, so in-sample), and
at $K = 16$ on the two holdout fixtures, where it is level at $K = 8$. At
$K = 4$ on those holdouts it does not: it ties on one and trails on the other.
On the hardware-derived holdout it is lower at $K = 16$
(Section~\ref{sec:eval}). That is
reported here and in the table rather than left to the appendix, because a
method whose advantage appears only above a budget threshold has a threshold,
and the reader is entitled to it.

\paragraph{What was not reached.}
The output names the representatives the search never judged and the number of
placements they stand for. They are \texttt{excluded\_by\_scope}: removed by a
budget, which is a property of the search, not of the hardware. A certificate
of optimality, where the search can produce one, is explicitly bounded by the
declared candidate space and the evaluation model---it says that nothing
\emph{evaluated under this model} can do better, and not that nothing can.
